% ==================== CHƯƠNG 5 ====================
\chapter{XÂY DỰNG CHƯƠNG TRÌNH MINH HỌA}

\section{Thư viện và công cụ sử dụng}

Trong quá trình phát triển ứng dụng quản lý lớp học thêm (hệ thống ``Lớp Mình''), nhóm chúng em đã sử dụng các ngôn ngữ lập trình, framework, thư viện, công cụ phát triển và cơ sở dữ liệu phù hợp để đảm bảo tính hiệu quả, bảo mật và dễ bảo trì. Các công cụ được chọn dựa trên yêu cầu dự án, bao gồm frontend xây dựng trên nền tảng Web hiện đại (HTML5, CSS3, JavaScript ES6+, Feather Icons), backend dựa trên Python với framework FastAPI và máy chủ Uvicorn, cùng hệ quản trị cơ sở dữ liệu quan hệ mạnh mẽ PostgreSQL (v14+). Dưới đây là danh sách chi tiết các thư viện và công cụ được sử dụng, được trình bày dưới dạng bảng.

\setlength{\tabcolsep}{3pt}
\renewcommand{\arraystretch}{1.25}
\begin{longtable}{|>{\raggedright\arraybackslash}p{3.8cm}|>{\raggedright\arraybackslash}p{5.4cm}|>{\raggedright\arraybackslash}p{6.4cm}|}
    \caption{Thư viện và công cụ sử dụng}\label{tab:c5-thu-vien-cong-cu}\\
    \hline\rowcolor{blue!11}
    \multicolumn{1}{|c|}{\textbf{Mục đích}} & \multicolumn{1}{c|}{\textbf{Công cụ}} & \multicolumn{1}{c|}{\textbf{Địa chỉ URL}} \\ \hline
    \endfirsthead

    \multicolumn{3}{r}{\textit{Tiếp bảng 5.1: Thư viện và công cụ sử dụng}}\\
    \hline\rowcolor{blue!11}
    \multicolumn{1}{|c|}{\textbf{Mục đích}} & \multicolumn{1}{c|}{\textbf{Công cụ}} & \multicolumn{1}{c|}{\textbf{Địa chỉ URL}} \\ \hline
    \endhead

    \hline
    \multicolumn{3}{r}{\footnotesize\textit{Còn tiếp trang sau}}\\
    \endfoot

    \hline
    \endlastfoot

    Môi trường & Python (v3.11+) \newline Node.js (v20+) & \url{https://www.python.org/} \newline \url{https://nodejs.org/} \\ \hline
    Ngôn ngữ lập trình & Python (v3.11.6) \newline JavaScript (ES6+ / CommonJS) \newline HTML5 / CSS3 \newline SQL (PostgreSQL Dialect) & \url{https://www.python.org/} \newline \url{https://developer.mozilla.org/en-US/docs/Web/JavaScript} \newline \url{https://developer.mozilla.org/} \newline \url{https://www.postgresql.org/docs/} \\ \hline
    Framework frontend & HTML5 Semantic \& Vanilla CSS3 \newline Hệ thống CSS Variables, Flexbox, Grid & \url{https://developer.mozilla.org/} \newline \url{https://developer.mozilla.org/en-US/docs/Web/CSS} \\ \hline
    Build tool \& Server frontend & Python HTTP Server (http.server v3.11+) & \url{https://docs.python.org/3/library/http.server.html} \\ \hline
    CSS framework & Custom Modular CSS Architecture (variables, base, components, dashboard) & \url{https://developer.mozilla.org/en-US/docs/Web/CSS} \\ \hline
    UI icons library & Feather Icons (v4.29.0) & \url{https://feathericons.com/} \\ \hline
    Charts \& Visualization & HTML5 Canvas \& Dynamic SVG & \url{https://developer.mozilla.org/en-US/docs/Web/API/Canvas_API} \\ \hline
    Routing \& Access control & Client-side Router \& Auth Guard (auth.js) & \url{https://developer.mozilla.org/en-US/docs/Web/API/History_API} \\ \hline
    Notification library & Custom Modal \& Toast Notifications (notifications.js) & \url{https://developer.mozilla.org/en-US/docs/Web/API/Notifications_API} \\ \hline
    HTTP client & Fetch API (Native Browser Web API) & \url{https://developer.mozilla.org/en-US/docs/Web/API/Fetch_API} \\ \hline
    Typography / Phông chữ & Google Fonts (Inter font family) & \url{https://fonts.google.com/specimen/Inter} \\ \hline
    Image upload \& File handling & FormData API \& File API (Native Web API) & \url{https://developer.mozilla.org/en-US/docs/Web/API/FormData} \\ \hline
    Client-side storage & Web Storage API (LocalStorage) & \url{https://developer.mozilla.org/en-US/docs/Web/API/Window/localStorage} \\ \hline
    Environment variables & Python os \& Dotenv configuration (.env) & \url{https://fastapi.tiangolo.com/advanced/settings/} \\ \hline
    Framework backend & FastAPI (v0.115.0+) & \url{https://fastapi.tiangolo.com/} \\ \hline
    Web server backend & Uvicorn [standard] (v0.30.0+) & \url{https://www.uvicorn.org/} \\ \hline
    Trình kết nối CSDL & Psycopg 3 (psycopg[binary] v3.2.0+) & \url{https://www.psycopg.org/} \\ \hline
    ORM \& Data Validation & Pydantic (v2.x) & \url{https://docs.pydantic.dev/} \\ \hline
    Authentication \& Security & Hashlib \& Secrets (scrypt KDF, Salt 16-byte, SHA-256) & \url{https://docs.python.org/3/library/hashlib.html} \\ \hline
    Session management & Bearer Session Token (SHA-256 token hash) & \url{https://datatracker.ietf.org/doc/html/rfc6750} \\ \hline
    File upload backend & Python-multipart (v0.0.12+) & \url{https://github.com/Kludex/python-multipart} \\ \hline
    Hỗ trợ \& Liên hệ trung tâm & Hotline 028 3822 8888 \& Trung tâm hỗ trợ 3 vai trò & \url{https://lopminh.edu.vn/} \\ \hline
    Caching \& Session & PostgreSQL Sessions (auth\_sessions) \& LocalStorage & \url{https://www.postgresql.org/} \\ \hline
    Unique ID generator & UUID v4 (Python uuid module) & \url{https://docs.python.org/3/library/uuid.html} \\ \hline
    CORS handling & FastAPI CORSMiddleware & \url{https://fastapi.tiangolo.com/tutorial/cors/} \\ \hline
    Rate limiting \& Brute-force & Lockout sau 3 lần sai / 15 phút (FastAPI \& Database) & \url{https://fastapi.tiangolo.com/} \\ \hline
    API documentation & Swagger UI / OpenAPI 3.0 \newline ReDoc & \url{https://swagger.io/} \newline \url{https://redocly.com/} \\ \hline
    Development server & Uvicorn Auto-reload (Backend) \newline Python HTTP Server (Frontend) & \url{https://www.uvicorn.org/} \newline \url{https://docs.python.org/3/library/http.server.html} \\ \hline
    Kiểm thử tự động & Python Unittest (API contract, DB schema, security) \newline Node.js Test Runner (node --test) & \url{https://docs.python.org/3/library/unittest.html} \newline \url{https://nodejs.org/api/test.html} \\ \hline
    Cơ sở dữ liệu chính & PostgreSQL (v14+) & \url{https://www.postgresql.org/} \\ \hline
    CSDL cấu trúc linh hoạt & PostgreSQL JSONB & \url{https://www.postgresql.org/docs/current/datatype-json.html} \\ \hline
    Kịch bản điều phối & PowerShell (run\_project.ps1) & \url{https://learn.microsoft.com/powershell/} \\ \hline
    IDE lập trình & Visual Studio Code (v1.90+) & \url{https://code.visualstudio.com/} \\ \hline
    Version control & Git (v2.56.0.windows.4) & \url{https://git-scm.com/} \\ \hline
    Remote repository & GitHub & \url{https://github.com/} \\ \hline
    Package manager & PIP (Python Package Installer v23.2+) & \url{https://pip.pypa.io/} \\ \hline
    Công cụ soạn thảo báo cáo & XeLaTeX / LaTeX (TeX Live) \newline Microsoft Word & \url{https://www.latex-project.org/} \newline \url{https://www.microsoft.com/microsoft-365/word} \\ \hline
\end{longtable}

\section{Kết quả chương trình minh họa}
Sau quá trình phân tích, thiết kế và hiện thực, nhóm đã xây dựng được phiên bản chương trình minh họa dưới dạng ứng dụng web mang tên \textit{Lớp Mình}, phục vụ công tác quản lý lớp học thêm. Hệ thống được tổ chức theo ba nhóm người dùng là giáo viên/gia sư, học sinh và phụ huynh. Người dùng đăng nhập bằng tài khoản được cấp và được chuyển đến không gian làm việc tương ứng với vai trò. Việc phân chia quyền truy cập giúp các nhóm người dùng tiếp cận những thông tin và thao tác phù hợp với trách nhiệm của mình.

Đối với giáo viên, chương trình đã xây dựng các nhóm chức năng quản lý tài khoản, lớp học, danh sách học sinh, thời khóa biểu và buổi học; thực hiện điểm danh; quản lý bài tập, bài kiểm tra, kết quả học tập và khoản thu học phí. Giao diện hỗ trợ lựa chọn lớp và buổi học, ghi nhận trạng thái chuyên cần, theo dõi kết quả đánh giá, lập khoản thu và cập nhật thông tin thanh toán. Đối với học sinh, chương trình cung cấp các màn hình tổng quan, lịch học, lịch sử điểm danh, kết quả học tập, bài tập, thông báo và hồ sơ cá nhân. Đối với phụ huynh, hệ thống bố trí các chức năng tra cứu lịch học, chuyên cần, kết quả học tập và tình trạng học phí của học sinh đã liên kết. Thông tin học phí được theo dõi trên hệ thống, còn giao dịch tiền mặt hoặc chuyển khoản được thực hiện bên ngoài ứng dụng và do giáo viên ghi nhận.

Về mặt kỹ thuật, phần giao diện được xây dựng bằng HTML, CSS và JavaScript, với các tệp xử lý được tổ chức thành những mô-đun dành cho từng vai trò. Phần máy chủ sử dụng ngôn ngữ Python và framework FastAPI để tiếp nhận, xử lý yêu cầu và cung cấp dịch vụ thông qua các API. Dữ liệu nghiệp vụ được tổ chức trong hệ quản trị cơ sở dữ liệu PostgreSQL. Trong bản mã nguồn khảo sát, tệp định nghĩa dịch vụ máy chủ có \textbf{64 khai báo API endpoint}, còn tệp thiết kế cơ sở dữ liệu có \textbf{22 bảng}. Bên cạnh đó, dự án có \textbf{27 trang HTML} theo ba vai trò (11 trang giáo viên, 8 trang học sinh và 8 trang phụ huynh), cùng một trang điều khoản dùng chung trong thư mục \texttt{pages}. Đây là kết quả thống kê cấu trúc mã nguồn, không đồng nghĩa mọi luồng nghiệp vụ đều đã được xác nhận hoạt động đầy đủ trên dữ liệu thực tế.

Mã nguồn chương trình được tổ chức thành các nhóm thư mục chính gồm \texttt{assets} (giao diện và mã JavaScript), \texttt{pages} (các trang theo vai trò), \texttt{backend} (dịch vụ và API), \texttt{database} (lược đồ và dữ liệu khởi tạo) và \texttt{tests} (các kịch bản kiểm tra). Nhóm lưu trữ và quản lý phiên bản mã nguồn trên GitHub tại địa chỉ \url{https://github.com/thichnguyettinh/Project_CNPM}. Cách phân chia này tạo điều kiện cho các thành viên phát triển, kiểm tra và tích hợp những phần chức năng được phân công.

Để đánh giá quy mô hiện thực, nhóm thống kê trực tiếp các tệp mã nguồn và kịch bản triển khai trong phiên bản dự án được khảo sát. Phạm vi thống kê bao gồm các tệp Python, JavaScript, HTML, CSS, SQL và PowerShell; không tính thư mục kiểm thử, tài liệu LaTeX, hình ảnh và các tệp nhị phân. Kết quả được trình bày tại Bảng~\ref{tab:c5-thong-ke-ma-nguon}.

\begin{table}[htbp]
    \centering
    \caption{Thống kê quy mô mã nguồn chương trình minh họa}
    \label{tab:c5-thong-ke-ma-nguon}
    \renewcommand{\arraystretch}{1.15}
    \begin{tabular}{|l|r|r|}
        \hline
        \textbf{Loại tệp} & \textbf{Số tệp} & \textbf{Số dòng} \\ \hline
        JavaScript (\texttt{.js}) & 30 & 7.395 \\ \hline
        HTML (\texttt{.html}) & 32 & 7.339 \\ \hline
        Python (\texttt{.py}) & 10 & 4.085 \\ \hline
        CSS (\texttt{.css}) & 19 & 3.928 \\ \hline
        SQL (\texttt{.sql}) & 7 & 861 \\ \hline
        PowerShell (\texttt{.ps1}) & 1 & 53 \\ \hline
        \textbf{Tổng cộng} & \textbf{99} & \textbf{23.661} \\ \hline
    \end{tabular}
\end{table}

Theo Bảng~\ref{tab:c5-thong-ke-ma-nguon}, dự án có tổng cộng \textbf{99 tệp} thuộc phạm vi thống kê, tương ứng \textbf{23.661 dòng mã nguồn vật lý}. Trong đó, JavaScript chiếm 7.395 dòng, HTML có 7.339 dòng, Python có 4.085 dòng và CSS có 3.928 dòng. Các tệp SQL và PowerShell lần lượt có 861 và 53 dòng. Số dòng được tính theo dòng vật lý của tệp văn bản, \textit{bao gồm cả dòng trống và dòng chú thích}; do đó đây là số liệu về quy mô tệp, không phải số dòng lệnh thực thi hay kết quả đo theo tiêu chí của công cụ \texttt{cloc}. Thống kê cho thấy phần giao diện và xử lý phía trình duyệt chiếm tỷ trọng đáng kể trong mã nguồn, phù hợp với đặc điểm một ứng dụng web hỗ trợ nhiều vai trò và nhiều màn hình nghiệp vụ.

Ngoài các chức năng nghiệp vụ, mã nguồn còn thể hiện việc tổ chức xác thực và quản lý phiên đăng nhập, phân quyền truy cập và lưu trữ thông tin mật khẩu theo dạng băm. Hệ thống có các giao diện thông báo trên website, hồ sơ cá nhân, đổi mật khẩu và trang hướng dẫn sử dụng. Các tệp trong thư mục \texttt{tests} hỗ trợ kiểm tra một số thành phần và hợp đồng API. Tuy nhiên, theo trạng thái được ghi nhận trong tài liệu mã nguồn, một số màn hình vẫn sử dụng dữ liệu minh họa hoặc chưa kết nối hoàn chỉnh với dịch vụ phía máy chủ; đồng thời các quy trình chưa được xác minh đầu cuối trên PostgreSQL cho mọi vai trò. Vì vậy, kết quả của giai đoạn này được xác định là \textbf{phiên bản chương trình minh họa đã có cấu trúc và các thành phần chức năng chính}, làm cơ sở cho việc tiếp tục kiểm thử, hoàn thiện tích hợp và triển khai ở các giai đoạn tiếp theo.

\setcounter{secnumdepth}{3}
\section{Giao diện minh hoạ các chức năng của chương trình}
\subsection{Giao diện Đăng nhập}

Đây là giao diện cho phép giáo viên, học sinh và phụ huynh truy cập vào hệ thống thông qua tên đăng nhập (hoặc email) và mật khẩu đã được cấp. Sau khi xác thực thành công, hệ thống tự động chuyển người dùng đến giao diện tương ứng với vai trò của mình.

Hệ thống không hỗ trợ đăng ký tài khoản công khai. Tài khoản học sinh và phụ huynh do giáo viên trực tiếp tạo và cấp phát. Người dùng chưa có tài khoản hoặc cần hỗ trợ mật khẩu có thể liên hệ giáo viên.

\begin{figure}[H]
  \centering
  \includegraphics[page=1,width=0.94\linewidth,height=0.49\textheight,keepaspectratio]{C5_AnhGiaoDien_28Trang.pdf}
  \caption{Giao diện đăng nhập}
  \label{fig:c5-01}
\end{figure}

\subsection{Giao diện Tổng quan}

\subsubsection{Giáo viên}

\begin{figure}[H]
  \centering
  \includegraphics[page=2,width=0.94\linewidth,height=0.49\textheight,keepaspectratio]{C5_AnhGiaoDien_28Trang.pdf}
  \caption{Giao diện Tổng quan của không gian giáo viên}
  \label{fig:c5-02}
\end{figure}

\subsubsection{Học sinh}

\begin{figure}[H]
  \centering
  \includegraphics[page=3,width=0.94\linewidth,height=0.49\textheight,keepaspectratio]{C5_AnhGiaoDien_28Trang.pdf}
  \caption{Giao diện Tổng quan của không gian học sinh}
  \label{fig:c5-03}
\end{figure}

\subsubsection{Phụ huynh}

\begin{figure}[H]
  \centering
  \includegraphics[page=4,width=0.94\linewidth,height=0.49\textheight,keepaspectratio]{C5_AnhGiaoDien_28Trang.pdf}
  \caption{Giao diện Tổng quan của không gian phụ huynh}
  \label{fig:c5-04}
\end{figure}

\subsection{Giao diện Lớp học}

\begin{figure}[H]
  \centering
  \includegraphics[page=5,width=0.94\linewidth,height=0.49\textheight,keepaspectratio]{C5_AnhGiaoDien_28Trang.pdf}
  \caption{Giao diện lớp học}
  \label{fig:c5-05}
\end{figure}

\subsection{Giao diện Tài khoản}

\begin{figure}[H]
  \centering
  \includegraphics[page=6,width=0.94\linewidth,height=0.49\textheight,keepaspectratio]{C5_AnhGiaoDien_28Trang.pdf}
  \caption{Giao diện quản lý tài khoản}
  \label{fig:c5-06}
\end{figure}

\subsection{Giao diện Lịch học}

\subsubsection{Giáo viên}

\begin{figure}[H]
  \centering
  \includegraphics[page=7,width=0.94\linewidth,height=0.49\textheight,keepaspectratio]{C5_AnhGiaoDien_28Trang.pdf}
  \caption{Giao diện lịch học của không gian giáo viên}
  \label{fig:c5-07}
\end{figure}

\subsubsection{Học sinh}

\begin{figure}[H]
  \centering
  \includegraphics[page=8,width=0.94\linewidth,height=0.49\textheight,keepaspectratio]{C5_AnhGiaoDien_28Trang.pdf}
  \caption{Giao diện Lịch học của học sinh}
  \label{fig:c5-08}
\end{figure}

\subsubsection{Phụ huynh}

\begin{figure}[H]
  \centering
  \includegraphics[page=9,width=0.94\linewidth,height=0.49\textheight,keepaspectratio]{C5_AnhGiaoDien_28Trang.pdf}
  \caption{Giao diện lịch học của không gian phụ huynh}
  \label{fig:c5-09}
\end{figure}

\subsection{Giao diện Điểm danh}

\subsubsection{Giáo viên}

\begin{figure}[H]
  \centering
  \includegraphics[page=10,width=0.94\linewidth,height=0.49\textheight,keepaspectratio]{C5_AnhGiaoDien_28Trang.pdf}
  \caption{Giao diện điểm danh của giáo viên}
  \label{fig:c5-10}
\end{figure}

\subsubsection{Học sinh}

\begin{figure}[H]
  \centering
  \includegraphics[page=11,width=0.94\linewidth,height=0.49\textheight,keepaspectratio]{C5_AnhGiaoDien_28Trang.pdf}
  \caption{Giao diện điểm danh của học sinh}
  \label{fig:c5-11}
\end{figure}

\subsubsection{Phụ huynh}

\begin{figure}[H]
  \centering
  \includegraphics[page=12,width=0.94\linewidth,height=0.49\textheight,keepaspectratio]{C5_AnhGiaoDien_28Trang.pdf}
  \caption{Giao diện điểm danh của không gian phụ huynh}
  \label{fig:c5-12}
\end{figure}

\subsection{Giao diện Điểm số}

\subsubsection{Giáo viên}

\begin{figure}[H]
  \centering
  \includegraphics[page=13,width=0.94\linewidth,height=0.49\textheight,keepaspectratio]{C5_AnhGiaoDien_28Trang.pdf}
  \caption{Giao diện quản lý điểm số của giáo viên}
  \label{fig:c5-13}
\end{figure}

\subsubsection{Giao diện Kết quả học tập và Bài tập (Không gian của Học sinh)}

\begin{figure}[H]
  \centering
  \includegraphics[page=14,width=0.94\linewidth,height=0.49\textheight,keepaspectratio]{C5_AnhGiaoDien_28Trang.pdf}
  \caption{Giao diện kết quả học tập của học sinh}
  \label{fig:c5-14}
\end{figure}

\begin{figure}[H]
  \centering
  \includegraphics[page=15,width=0.94\linewidth,height=0.49\textheight,keepaspectratio]{C5_AnhGiaoDien_28Trang.pdf}
  \caption{Giao diện bài tập của học sinh}
  \label{fig:c5-15}
\end{figure}

\subsubsection{Giao diện Kết quả học tập (Không gian của phụ huynh)}

\begin{figure}[H]
  \centering
  \includegraphics[page=16,width=0.94\linewidth,height=0.49\textheight,keepaspectratio]{C5_AnhGiaoDien_28Trang.pdf}
  \caption{Giao diện Kết quả học tập của không gian phụ huynh}
  \label{fig:c5-16}
\end{figure}

\subsection{Giao diện Học phí}

\subsubsection{Giáo viên}

\begin{figure}[H]
  \centering
  \includegraphics[page=17,width=0.94\linewidth,height=0.49\textheight,keepaspectratio]{C5_AnhGiaoDien_28Trang.pdf}
  \caption{Giao diện quản lý học phí của không gian giáo viên}
  \label{fig:c5-17}
\end{figure}

\subsubsection{Phụ huynh}

\begin{figure}[H]
  \centering
  \includegraphics[page=18,width=0.94\linewidth,height=0.49\textheight,keepaspectratio]{C5_AnhGiaoDien_28Trang.pdf}
  \caption{Giao diện chi tiết học phí của không gian phụ huynh}
  \label{fig:c5-18}
\end{figure}

\subsection{Giao diện Thông báo}

\subsubsection{Giáo viên}

\begin{figure}[H]
  \centering
  \includegraphics[page=19,width=0.94\linewidth,height=0.49\textheight,keepaspectratio]{C5_AnhGiaoDien_28Trang.pdf}
  \caption{Giao diện thông báo của không gian giáo viên}
  \label{fig:c5-19}
\end{figure}

\subsubsection{Học sinh}

\begin{figure}[H]
  \centering
  \includegraphics[page=20,width=0.94\linewidth,height=0.49\textheight,keepaspectratio]{C5_AnhGiaoDien_28Trang.pdf}
  \caption{Giao diện thông báo của học sinh}
  \label{fig:c5-20}
\end{figure}

\subsubsection{Phụ huynh}

\begin{figure}[H]
  \centering
  \includegraphics[page=21,width=0.94\linewidth,height=0.49\textheight,keepaspectratio]{C5_AnhGiaoDien_28Trang.pdf}
  \caption{Giao diện thông báo của phụ huynh}
  \label{fig:c5-21}
\end{figure}

\subsection{Giao diện Hồ sơ cá nhân}

\subsubsection{Giáo viên}

\begin{figure}[H]
  \centering
  \includegraphics[page=22,width=0.94\linewidth,height=0.49\textheight,keepaspectratio]{C5_AnhGiaoDien_28Trang.pdf}
  \caption{Giao diện hồ sơ của giáo viên}
  \label{fig:c5-22}
\end{figure}

\subsubsection{Học sinh}

\begin{figure}[H]
  \centering
  \includegraphics[page=23,width=0.94\linewidth,height=0.49\textheight,keepaspectratio]{C5_AnhGiaoDien_28Trang.pdf}
  \caption{Giao diện hồ sơ của học sinh}
  \label{fig:c5-23}
\end{figure}

\subsubsection{Phụ huynh}

\begin{figure}[H]
  \centering
  \includegraphics[page=24,width=0.94\linewidth,height=0.49\textheight,keepaspectratio]{C5_AnhGiaoDien_28Trang.pdf}
  \caption{Giao diện hồ sơ của phụ huynh}
  \label{fig:c5-24}
\end{figure}

\subsection{Giao diện Đổi mật khẩu}

\begin{figure}[H]
  \centering
  \includegraphics[page=25,width=0.94\linewidth,height=0.49\textheight,keepaspectratio]{C5_AnhGiaoDien_28Trang.pdf}
  \caption{Giao diện đổi mật khẩu}
  \label{fig:c5-25}
\end{figure}

\subsection{Giao diện Trung tâm hướng dẫn}

\subsubsection{Giáo viên}

\begin{figure}[H]
  \centering
  \includegraphics[page=26,width=0.94\linewidth,height=0.49\textheight,keepaspectratio]{C5_AnhGiaoDien_28Trang.pdf}
  \caption{Giao diện trung tâm hướng dẫn của giáo viên}
  \label{fig:c5-26}
\end{figure}

\subsubsection{Học sinh}

\begin{figure}[H]
  \centering
  \includegraphics[page=27,width=0.94\linewidth,height=0.49\textheight,keepaspectratio]{C5_AnhGiaoDien_28Trang.pdf}
  \caption{Giao diện trung tâm hướng dẫn của học sinh}
  \label{fig:c5-27}
\end{figure}

\subsubsection{Phụ huynh}

\begin{figure}[H]
  \centering
  \includegraphics[page=28,width=0.94\linewidth,height=0.49\textheight,keepaspectratio]{C5_AnhGiaoDien_28Trang.pdf}
  \caption{Giao diện trung tâm hướng dẫn của phụ huynh}
  \label{fig:c5-28}
\end{figure}
